% !TeX program = xelatex
\documentclass[12pt,a4paper]{article}

% ==================== 基础宏包 ====================
\usepackage[UTF8,fontset=fandol]{ctex} % 使用 TeX Live 自带字体，便于跨环境编译。
\usepackage[margin=2.2cm]{geometry}
\usepackage{amsmath,amssymb,amsthm,mathtools}
\usepackage{array}
\usepackage{booktabs}
\usepackage{enumitem}
\usepackage{xcolor}
\usepackage{graphicx}
\usepackage{fancyhdr}
\usepackage{lastpage}
\usepackage{hyperref}
\usepackage[most]{tcolorbox}
\usepackage{titlesec}
\usepackage{tabularx}

\hypersetup{
  hidelinks
}

% ==================== 颜色定义 ====================
\definecolor{mainblue}{RGB}{34,72,112}
\definecolor{lightblue}{RGB}{241,246,252}
\definecolor{linegray}{RGB}{210,218,228}
\definecolor{textgray}{RGB}{90,98,108}

% ==================== 页面设置 ====================
\setlength{\parindent}{2em}
\setlength{\parskip}{0.35em}
\setlength{\headheight}{25pt}
\linespread{1.2}

\pagestyle{fancy}
\fancyhf{}
\lhead{\textcolor{textgray}{\small 课程作业}}
\rhead{\textcolor{textgray}{\small \thepage/\pageref*{LastPage}}}
\cfoot{}
\renewcommand{\headrulewidth}{0.4pt}
\renewcommand{\headrule}{\hbox to\headwidth{\color{linegray}\leaders\hrule height \headrulewidth\hfill}}

% ==================== 标题样式 ====================
\titleformat{\section}
  {\Large\bfseries\color{mainblue}}
  {\thesection}
  {0.6em}
  {}

\titleformat{\subsection}
  {\large\bfseries\color{mainblue}}
  {\thesubsection}
  {0.6em}
  {}
\renewcommand{\thesubsection}{(\alph{subsection})}

\titlespacing*{\section}{0pt}{1.2em}{0.6em}
\titlespacing*{\subsection}{0pt}{0.8em}{0.4em}

\setlist[enumerate,1]{label=(\alph*),leftmargin=2em,
  topsep=0.25em,itemsep=0.1em,parsep=0pt,partopsep=0pt}

% ==================== 自定义环境 ====================
\newtcolorbox{infobox}{
  colback=lightblue,
  colframe=mainblue,
  boxrule=0.6pt,
  arc=2mm,
  left=1.2em,
  right=1.2em,
  top=0.8em,
  bottom=0.8em
}

\newtcolorbox{answerbox}{
  colback=white,
  colframe=linegray,
  boxrule=0.5pt,
  arc=2mm,
  left=1em,
  right=1em,
  top=0.55em,
  bottom=0.55em,
  before skip=0.35em,
  after skip=0.35em,
  breakable
}

\newcommand{\course}{随机过程}
\newcommand{\homeworktitle}{第9次作业}
\newcommand{\studentname}{倪伟丰}
\newcommand{\studentid}{2024110089}
\newcommand{\classname}{24 大数据 2 班}
\newcommand{\teacher}{郭珊珊}
\newcommand{\duedate}{} % 在此填写提交日期

% ==================== 正文开始 ====================
\begin{document}

\begin{center}
  {\Huge\bfseries\color{mainblue}\homeworktitle}\\[0.6em]
  {\Large\course}\\[1.2em]
  {\color{linegray}\rule{0.9\textwidth}{0.8pt}}
\end{center}

\begin{infobox}
\begin{tabularx}{\linewidth}{>{\bfseries}p{2.2cm} X >{\bfseries}p{2.2cm} X}
姓名 & \studentname & 学号 & \studentid \\
班级 & \classname & 教师 & \teacher \\
提交日期 & 11月10日 &  &  \\
\end{tabularx}
\end{infobox}

% 题目来源：同目录 assignment9-1234&1247-1.docx。
% 版式参考：../assignment7/main.tex。每个 answerbox 内可直接填写解答。


\section{Problem 78}

A store opens at 8~A.M. From 8 until 10~A.M. customers arrive at a Poisson
rate of four an hour. Between 10~A.M. and 12~P.M. they arrive at a Poisson
rate of eight an hour. From 12~P.M. to 2~P.M. the arrival rate increases
steadily from eight per hour at 12~P.M. to ten per hour at 2~P.M.;
and from 2 to 5~P.M. the arrival rate drops steadily from ten per hour
at 2~P.M. to four per hour at 5~P.M. Determine the probability distribution
of the number of customers that enter the store on a given day.

\begin{answerbox}
\textbf{解：}

$t$ 表示当天的钟点，以小时为单位，并令 $N(t)$ 表示从早上
8 点到时刻 $t$ 累计进入商店的顾客人数，$N(8)=0$。
顾客到达构成非齐次泊松过程，其强度（到达速率）函数为

$$
\lambda(t)=
\begin{cases}
4, & 8\le t<10,\\
8, & 10\le t<12,\\
8+(t-12)=t-4, & 12\le t<14,\\
10-2(t-14)=38-2t, & 14\le t\le17.
\end{cases}
$$

均值函数为
$$
m(t)=\int_8^t\lambda(u)\,du
=\begin{cases}
4(t-8), & 8\le t<10,\\[2pt]
8+8(t-10), & 10\le t<12,\\[2pt]
24+8(t-12)+\dfrac{(t-12)^2}{2}, & 12\le t<14,\\[6pt]
42+10(t-14)-(t-14)^2, & 14\le t\le17.
\end{cases}
$$

将 $t=17$ 代入上式，得到 $m(17)=63$。
由非齐次泊松过程的计数分布，全天进入商店的顾客总数满足
$$
N(17)\sim\operatorname{Poisson}(63),
\qquad
P\{N(17)=n\}=e^{-63}\frac{63^n}{n!},\quad n=0,1,2,\ldots.
$$

\end{answerbox}

\clearpage
\section{Problem 84}

Let $X_1, X_2, \ldots$ be independent and identically distributed nonnegative
continuous random variables having density function $f(x)$. We say that a
record occurs at time $n$ if $X_n$ is larger than each of the previous values
$X_1, \ldots, X_{n-1}$. (A record automatically occurs at time~1.)
If a record occurs at time $n$, then $X_n$ is called the record value.
In other words, a record occurs whenever a new high is reached, and that
new high is called the record value. Let $N(t)$ denote the number of record
values that are less than or equal to $t$. Characterize the process
$\{N(t), t \ge 0\}$ when

\subsection{}
$f$ is an arbitrary continuous density function.

\begin{answerbox}
\textbf{解：}

记 $F$ 为 $X_i$ 的分布函数。这里 $t$ 是纪录值的上界。
先考虑 $F(t)<1$ 的情形。

第一条超过 $t$ 的纪录来自第一个满足 $X_i>t$ 的样本，其分布为
$X_1$ 在 $X_1>t$ 条件下的分布。由独立同分布性，该样本及后续
样本与此前不超过 $t$ 的纪录独立，所以阈值 $t$ 以上的纪录与
$t$ 以内的纪录独立，即 $\{N(t)\}$ 具有独立增量。

区间 $(t,t+h]$ 内至少出现一条纪录，当且仅当第一条超过 $t$ 的
纪录不超过 $t+h$，故当 $h\downarrow0$ 时，
$$
P\{N(t+h)-N(t)\ge1\}
=\frac{F(t+h)-F(t)}{1-F(t)}
=r(t)h+o(h),
\qquad r(t)=\frac{f(t)}{1-F(t)}.
$$
出现第一条纪录后，阈值升高，再出现一条的条件概率不超过上述
比值。因此，区间内出现两条及以上纪录的概率为 $O(h^2)=o(h)$，
恰好出现一条的概率为 $r(t)h+o(h)$。

又 $N(0)=0$，所以这是强度函数为 $r(t)$ 的非齐次泊松过程，且
$$
m(t)=\int_0^t r(u)\,du=-\ln(1-F(t)),
\qquad N(t)\sim\operatorname{Poisson}(m(t)).
$$
若存在有限的 $t$ 使 $F(t)=1$，则所有纪录值都不超过 $t$，
此时 $N(t)=\infty$（几乎必然）；上述有限参数的泊松分布只适用于 $F(t)<1$。

\end{answerbox}

\subsection{}

$f(x) = \lambda e^{-\lambda x}$.

\begin{answerbox}
\textbf{解：}

原始样本 $X_i$ 服从指数分布，故 $r(t)=f(t)/(1-F(t))=\lambda$。
因此，$\{N(t)\}$ 是速率为 $\lambda$ 的齐次泊松过程，且
$$
N(t) \sim \operatorname{Poisson}(\lambda t).
$$

\end{answerbox}


\clearpage
\section{Problem 85}

Let $X(t) = \sum_{i=1}^{N(t)} X_i$ where $X_i$, $i \ge 1$ are independent
and identically distributed with mean $E[X]$, and are independent of
$\{N(t), t \ge 0\}$, which is a Poisson process with rate $\lambda$.
For $s < t$, find

\subsection{}

$E[X(t) \mid X(s)]$;

\begin{answerbox}
\textbf{解：}

将时刻 $s$ 之后的增量记为
$$
X_s(t)=\sum_{i=N(s)+1}^{N(t)}X_i,
\qquad X(t)=X(s)+X_s(t).
$$
由泊松过程的独立增量性，以及 $X_i$ 彼此独立且与计数过程独立，
$X_s(t)$ 与 $X(s)$ 独立。
又 $N(t)-N(s)\sim\operatorname{Poisson}(\lambda(t-s))$；
给定新增项数为 $k$ 时，$X_s(t)$ 的均值为 $kE[X]$，所以
$$
E[X_s(t)]=E[N(t)-N(s)]E[X]=\lambda(t-s)E[X].
$$
于是
$$
E[X(t)\mid X(s)]
=X(s)+E[X_s(t)]
=X(s)+\lambda(t-s)E[X].
$$

\end{answerbox}

\subsection{}

$E[X(t) \mid N(s)]$;

\begin{answerbox}
\textbf{解：}

当 $N(s)=n$ 时，由独立增量性，
$$
E[N(t)\mid N(s)=n]
=n+E[N(t)-N(s)]
=n+\lambda(t-s).
$$
由于 $\{X_i\}$ 与计数过程独立，给定 $N(t)$ 和 $N(s)$ 后，
$X(t)$ 的条件均值为 $N(t)E[X]$。由条件期望的迭代法则，
$$
\begin{aligned}
E[X(t)\mid N(s)]
&=E\bigl[E[X(t)\mid N(t),N(s)]\mid N(s)\bigr]\\
&=E[N(t)\mid N(s)]E[X]\\
&=\bigl(N(s)+\lambda(t-s)\bigr)E[X].
\end{aligned}
$$


\end{answerbox}

\clearpage
\section{Problem 3}

Consider two machines that are maintained by a single repairman. Machine $i$
functions for an exponential time with rate $\mu_i$ before breaking down,
$i = 1, 2$. The repair times (for either machine) are exponential with rate
$\mu$. Can we analyze this as a birth and death process? If so, what are the
parameters? If not, how can we analyze it?

\begin{answerbox}
\textbf{解：}

假设各机器的工作时间与维修时间相互独立；维修工在有故障机器时
立即维修，且每次修完一台再修另一台。

令 $B(t)$ 表示时刻 $t$ 的故障机器数。当 $B(t)=1$ 时，若机器 1
正在维修、机器 2 正常，则下一次故障的速率为 $\mu_2$；若机器 2
正在维修、机器 1 正常，则下一次故障的速率为 $\mu_1$。
因此，当 $\mu_1\ne\mu_2$ 时，仅知道故障机器数无法确定下一次
故障的速率，还需知道正常机器的身份。因此，$\{B(t)\}$ 一般不具有
马尔可夫性，不能仅按故障机器数建立生灭过程。

可以扩大状态空间，建立连续时间马尔可夫链。定义状态如下：
\begin{center}
\begin{tabular}{cl}
\toprule
状态 & 机器情况 \\
\midrule
0 & 两台机器均正常 \\
1 & 机器 1 正在维修，机器 2 正常 \\
2 & 机器 2 正在维修，机器 1 正常 \\
3 & 两台机器均故障，正在维修机器 1 \\
4 & 两台机器均故障，正在维修机器 2 \\
\bottomrule
\end{tabular}
\end{center}
两台机器均故障时，必须区分正在维修哪一台，因为这决定了维修
完成后进入的状态。由指数分布的无记忆性，按状态 $0,1,2,3,4$
的顺序，生成矩阵为
$$
Q=
\begin{pmatrix}
-(\mu_1+\mu_2)&\mu_1&\mu_2&0&0\\
\mu&-(\mu+\mu_2)&0&\mu_2&0\\
\mu&0&-(\mu+\mu_1)&0&\mu_1\\
0&0&\mu&-\mu&0\\
0&\mu&0&0&-\mu
\end{pmatrix}.
$$

特殊地，若 $\mu_1=\mu_2=\nu$，则故障机器数 $\{B(t)\}$
构成状态空间为 $\{0,1,2\}$ 的生灭过程。以故障为出生、维修
完成为死亡，其出生率与死亡率分别为
$$
\begin{gathered}
b_0=2\nu,\qquad b_1=\nu,\qquad b_2=0,\\
d_0=0,\qquad d_1=d_2=\mu.
\end{gathered}
$$
只有一名维修工，因此两台均故障时的总维修率仍为 $\mu$。

\end{answerbox}

\clearpage
\section{Problem 5}

There are $N$ individuals in a population, some of whom have a certain
infection that spreads as follows. Contacts between two members of this
population occur in accordance with a Poisson process having rate $\lambda$.
When a contact occurs, it is equally likely to involve any of the
$\binom{N}{2}$ pairs of individuals in the population. If a contact involves
an infected and a noninfected individual, then with probability $p$ the
noninfected individual becomes infected. Once infected, an individual
remains infected throughout. Let $X(t)$ denote the number of infected
members of the population at time $t$.

\subsection{}

Is $\{X(t), t \ge 0\}$ a continuous-time Markov chain?

\begin{answerbox}
\textbf{解：}

是。状态空间为 $\{0,1,\ldots,N\}$。当 $X(t)=i$ 时，有 $i$ 名
感染者和 $N-i$ 名未感染者，因此一名感染者与一名未感染者构成的
配对共有 $i(N-i)$ 对。一次接触导致新增一名感染者的概率为
$$
r_i=p\frac{i(N-i)}{\binom{N}{2}}.
$$
题中的 $\lambda$ 是全体人群的总接触率。在状态 $i$ 保持不变期间，
对速率为 $\lambda$ 的接触泊松过程
按概率 $r_i$ 进行筛选，得到有效传播事件的速率
$$
b_i=\lambda r_i=\lambda p\frac{i(N-i)}{\binom{N}{2}}.
$$
若 $b_i>0$，等待下一次有效传播的时间服从
$\operatorname{Exp}(b_i)$，传播发生后进入状态 $i+1$。
若 $b_i=0$，则一直停留在状态 $i$。
上述等待时间和下一状态的条件分布仅依赖于当前感染人数 $i$，
与此前的接触及感染历史无关。故 $\{X(t),t\ge0\}$ 是连续时间
马尔可夫链。
\end{answerbox}

\subsection{}

Specify its type.

\begin{answerbox}
\textbf{解：}

每次有效传播使感染人数增加 1，且感染者不会恢复，因此这是一个
有限状态的纯生过程。出生率为 (a) 中的 $b_i$，死亡率为
$d_i=0$（$i=0,1,\ldots,N$）。
生成矩阵满足 $q_{i,i+1}=b_i$（$i<N$）、$q_{ii}=-b_i$，其余元素为 0。
由于 $b_0=b_N=0$，状态 0 和 $N$ 均为吸收态：没有感染者时
无法传播，所有人均感染后人数也无法再增加。
\end{answerbox}

\end{document}
